\subsection{The Adjoint Operator}

\begin{itemize}
    \item From the Hahn-Banach Theorem, we can see that for every bounded linear operator \( T: X \to Y  \) on a normed space \( X  \), we can make sense of and associate an adjoint operator \( T^{\times} \) of \( T  \). 
    \item This helps us solve equations in physics and other applications.
    \item For this section, we will look at some of its properties and how it is related to the Hilbert-adjoint operator \( T^{*} \) defined in earlier sections of this book. \textbf{Note:} When we refer to the Adjoint Operator of \( T  \), we are not referring to the Hilbert-adjoint Operator of \( T  \). These two concepts are distinct from each other. 
    \item To this end, let us first consider a bounded linear operator \( T : X \to Y  \), where \( X  \) and \( Y  \) are normed spaces.
\end{itemize}

Let \( g  \) be any bounded linear functional on \( Y  \). Since \( g  \) is defined on all of \( Y  \), we can set \( y = Tx  \) and obtain a functional \( f  \) such that 
\begin{equation}  f(x) = g(Tx) \ \ \forall x \in X. \tag{1}  \end{equation}
\begin{itemize}
    \item Notice that \( f  \) is linear since \( g  \) and \( T  \) are linear. 
    \item \( f  \) is bounded since \( g  \) and \( T  \) are both bounded. That is, 
        \[  | f(x) |  = | g (Tx) |  \leq \|g \| \| Tx \| \leq \|g \| \|T \| \|x \|.    \]
        Taking the supremum over all \( x \in X  \) where \( \|x \| = 1  \), we obtain
        \[  \|f\| \leq \|g \| \|T\|. \tag{2} \]
        Thus, (2) tells us that \( f \in X' \), where \( X' \) is the dual  space of \( X  \). 
    \item (1) tells us that \( f  \) is an operator defined from \( Y' \) into \( X' \). This is precisely the \textbf{adjoint operator} of \( T \) which will be defined more explicitly below. 
\end{itemize}

\begin{definition}[Adjoint Operator \( T^{\times} \)]
    Let \( T : X \to Y  \) be a bounded linear operator, where \( X  \) and \( Y  \) are normed spaces. Then the \textbf{adjoint operator} \( T^{\times}: Y' \to X' \) of \( T  \) is defined by
    \[  f(x) = (T^{\times}g)(x) = g(Tx) \ \ \forall g \in Y'. \]
\end{definition}

We will see in the next theorem that the adjoint operator \( T^{\times} \) of \( T  \) will have the same norm as the operator itself. To do this, we will need {\hyperref[4.3-3]{Theorem 4.3-3}} (a result that followed from the Hahn-Banach Theorem) to establish this property.

\subsection{Equivalent Norms}

\begin{theorem}[Norm Of The Adjoint Operator (4.5-2)]
    The adjoint operator \( T^{\times}  \) is linear and bounded, and 
    \[  \|T^{\times}\| = \|T\|. \tag{5}\]
\end{theorem}
\begin{proof}
The operator \( T^{\times} \) is linear since its domain \( Y' \) is a vector space. Thus, we obtain
\begin{align*}
    (T^{\times}(\alpha {g}_{1} + \beta {g}_{2}))(x) &= (\alpha {g}_{1} + \beta {g}_{2})(Tx) \\
                                                    &= \alpha {g}_{1} (Tx) + \beta {g}_{2}(Tx) \\
                                                    &= \alpha (T^{\times}{g}_{1})(x) + \beta (T^{\times} {g}_{2})(x).
\end{align*}
We will prove (5) from (4). That is, we will show that \( \|T^{\times}\| \leq \|T\| \) and \( \|T^{\times}\| \geq \|T\| \). For \( f = T^{\times} g  \), we can use (2) to write
\[  \|T^{\times} g \| = \|f\| \leq \|g\| \|T\|. \]
Taking the supremum over all \( g \in Y' \) of norm one, we obtain 
\[  \|T^{\times}\| \leq \|T\| \]
proving the first inequality. To get the other inequality, we will use Theorem 4.3-3. That is, for every \( {x}_{0} \neq 0  \) in \( X  \), there exists a \( {g}_{0} \in Y' \) such that 
\[  \|{g}_{0}\| = 1  \ \ \text{and} \ \ {g}_{0}(T {x}_{0}) = \|T {x}_{0}\|. \]
Notice, by definition, that \( {g}_{0}(T {x}_{0}) = (T^{\times} {g}_{0})({x}_{0}) \). Setting \( {f}_{0} = T^{\times}{g}_{0} \), we obtain
\begin{align*}
    \|T {x}_{0} \| = {g}_{0}(T {x}_{0}) &= {f}_{0}({x}_{0}) \\
                                    &\leq \|{f}_{0}({x}_{0})\| \\
                                        &\leq \|{f}_{0}\| \|{x}_{0}\| \\
                                        &= \|T^{\times} {g}_{0}\| \|{x}_{0}\| \\
                                        &\leq \|T^{\times}\| \|{g}_{0}\| \|{x}_{0}\|.
\end{align*}
Since \( \|{g}_{0}\| = 1  \), we further have
\[  \|T {x}_{0} \| \leq \|T^{\times}\| \|{x}_{0}\| \]
which holds for \( {x}_{0} = 0  \) since \( T 0 = 0  \). Also, we have
\[  \|T {x}_{0} \| \leq \|T \| \|{x}_{0}\| \]
where \( \|T\| \) is the least upper bound by definition and \( \|T^{\times}\|  \) is an upper bound of \( \|T {x}_{0}\| \). Hence, \( \|T^{\times}\| \geq \|T\| \) giving us the second inequality. Thus, we conclude that \( \|T^{\times}\| = \|T\| \).
\end{proof}

\begin{eg}[Adjoint of a Linear Operator in \( \R^n \)]
  Given a linear operator \( T: \R^{n} \to \R^{n} \), the matrix representation \( {T}_{E} = ({\tau}_{jk})  \), its adjoint is just the transpose of \( {T}_{E} \) where \( E = \{  {e}_{1}, \dots, {e}_{n} \}   \) is a basis for \( \R^{n} \) whose elements are arranged in an order that is fixed. Let \( x = ({\xi}_{1}, \dots, {\xi}_{n}) \) and \( y = ({\eta}_{1} ,\dots, {\eta}_{n}) \) be two column vectors in \( \R^{n} \). Let \( F = \{  {f}_{1}, \dots, {f}_{n} \}   \) be the dual basis for the dual of \( \R^{n} \). Given that every linear functional \( g   \) in the dual of \( \R^{n} \) has the representation, 
  \[  g = {\alpha}_{1} {f}_{1} + \cdots + {\alpha}_{n} {f}_{n},  \]
 the formula for such an adjoint is given by   
  \[  {T}_{E}^{\times} g (x) = g ({T}_{E}x) = \sum_{ k=1  }^{ n  } {\beta}_{k} {\xi}_{k} \ \ \text{where} \ \ {\beta}_{k} = \sum_{ k=1  }^{ n } {\tau}_{jk} {\alpha}_{j}. \]

  This tells us that for every linear operator \( T  \) from \( \R^{n} \) to \( \R^{n} \) with matrix representation \( {T}_{E} \), the adjoint operator \( T^{\times} \) is the transpose of \( {T}_{E} \). The same result holds if \( T  \) is a linear operator from \( \C^{n} \) to \( \C^{n} \). 
\end{eg}

Below are some properties of the adjoint operator:

\begin{itemize}
    \item For every \( S,T \in B(X,Y) \), we have \( (S+T)^{\times} = S^{\times} + T^{\times} \).
    \item \( (\alpha T)^{\times} = \alpha T^{\times} \). 
    \item Let \( X,Y,Z  \) be normed spaces and \( T \in B(X,Y) \) and \( S \in B(Y,Z) \). Then the adjoint operator of \( ST \) is
        \[  (ST)^{\times} = T^{\times} S^{\times}. \]
    \item \( (T^{\times})^{-1} = (T^{-1})^{\times} \).
\end{itemize}

\subsection{Relation Between Adjoint Operator and Hilbert-Adjoint Operator}

If \( X  \) and \( Y  \) are both Hilbert spaces, then we can express the Hilbert-adjoint operator \( T^{*} \) of a linear operator \( T  \) on a Hilbert space in terms of the adjoint operator \( T^{\times} \) of \( T  \).

Let \( X = {H}_{1} \) and \( Y = {H}_{2} \) where \( {H}_{1} \) and \( {H}_{2} \) are two Hilbert spaces. Consider \( T: {H}_{1} \to {H}_{2} \) and \( {T}_{\times}: {H}_{2}' \to {H}_{1}' \) and recall that the adjoint \( T^{\times} \) is defined by 
\begin{align*}
    T^{\times} g &= f \tag{a} \\
    g(Tx) &= f(x) \tag{b}
\end{align*}
where \( f \in {H}_{1}' \) and \( g \in {H}_{2}' \).  Now that we have two Hilbert spaces \( {H}_{1} \) and \( {H}_{2} \) to work with, we can express each linear functional \( f  \) and \( g  \) from \( {H}_{1} \) and \( {H}_{2} \), respectively, via the Riesz representation; that is, we can find \( {x}_{0} \in {H}_{1} \) and \( {y}_{0} \in {H}_{2} \) such that  
\begin{align*}
    f(x) &= \langle x , {x}_{0} \rangle \\
    g(y) &= \langle y , {y}_{0} \rangle.
\end{align*}

Furthermore, we see that \( {x}_{0} \) and \( {y}_{0}  \) are uniquely determined by \( f  \) and \( g  \) (using Theorem 3.8-1), respectively, thereby giving us the following operators \( {A}_{1} : {H}_{1}' \to {H}_{1} \) and \( {A}_{2} : {H}_{2}' \to {H}_{2} \) defined by \( {A}_{1} f = {x}_{0} \) and \( {A}_{2} g = {y}_{0} \). From the same theorem, \( {A}_{1} \) and \( {A}_{2} \) are both:
\begin{enumerate}
    \item[(a)] bijections;
    \item[(b)] isometries because \( \|{A}_{1}f\| = \|{x}_{0}\| = \|f\| \) and \( \|{A}_{2}g\| = \|{y}_{0}\| = \|g\| \);
    \item[(c)] conjugate linear.
\end{enumerate}
Composing \( {A}_{2}^{-1}, T^{\times},   \) and \( {A}_{1} \) together, we define
\[ {A}_{1} T^{\times} {A}_{2}^{-1} : {H}_{2} \to {H}_{1} \ \ \text{by} \ \ {A}_{1}T^{\times} {A}_{2}^{-1} {y}_{0} = {x}_{0}. \]
Our claim is that \( T^{*} = {A}_{1} T^{\times} {A}_{2}^{-1} \). Since we are able to take advantage of the Riesz representations of each linear functional from their respective duals, it follows that 
\[  \langle Tx , {y}_{0} \rangle = g(Tx) = f(x) = \langle x , {x}_{0} \rangle = \langle x  ,  ({A}_{1} T^{\times} {A}_{2} )  {y}_{0} \rangle. \]
Since \( \langle Tx  , {y}_{0} \rangle = \langle x  , T^{*} {y}_{0} \rangle \), we conclude that \( {A}_{1} T^{\times} {A}_{2}^{-1} = T^{*} \).

\subsection{Key Differences Between Adjoints and Hilbert-Adjoint}

\begin{itemize}
    \item The Hilbert-adjoint \( T^{*} \) is defined directly on the spaces \( X  \) and \( Y  \) whereas the adjoint \( T^{\times} \) is defined on the duals of \( X  \) and \( Y \), respectively.
    \item For \( \alpha \in \F  \), we have \( (\alpha T)^{\times} = \alpha T^{\times} \) whereas \( (\alpha T^{*}) = \overline{\alpha} T^{*} \).
    \item For finite dimensional Euclidean spaces, the adjoint \( T^{\times} \) represents the transpose the matrix representation of \( T  \) and the Hilbert-adjoint \( T^{*} \) represents the conjugate transpose of the same matrix. 
\end{itemize}
